\subsection{Batzner et al.\ (2022): NequIP E(3)-equivariant GNN MLIP}
\label{sec:appendix-example-batzner2022}

\paragraph{Q1 -- Bibliographic identification.}
Batzner, Musaelian, Sun, Geiger, Mailoa, Kornbluth, Molinari, Smidt,
and Kozinsky introduce NequIP, an $E(3)$-equivariant graph
neural-network interatomic potential. Published in
\emph{Nature Communications} 13, 2453 (2022); arXiv:2101.03164.
The paper benchmarks NequIP across MD17/rMD17 small molecules,
Molecules@CCSD/CCSD(T), liquid water + ice phases, formate
decomposition on Cu(110), Li$_4$P$_2$O$_7$ amorphous glass, and
the LiPS superionic conductor.
\lstinputlisting[language=turtle]{artifacts/kg/listings/batzner2022/q01-bibliographic.ttl}

\paragraph{Q2 -- Material system.}
The canonical encoding pass takes the bulk water and ice phases
benchmark as the headline extended-system application: 64 H$_2$O
molecules in the liquid phase, 96 H$_2$O molecules in three ice Ih
phases (1~bar/273~K, 1~bar/330~K, 2.13~kbar/238~K).
\lstinputlisting[language=turtle]{artifacts/kg/listings/batzner2022/q02-material.ttl}

\paragraph{Q3 -- Reference calculation method.}
Reference data are produced from DFT calculations.
\lstinputlisting[language=turtle]{artifacts/kg/listings/batzner2022/q03-reference-method-type.ttl}

\paragraph{Q4 -- Reference settings.}
The water + ices reference data are computed at the PBE0-TS level
of theory (hybrid PBE0 with Tkatchenko--Scheffler dispersion),
generated from a mix of classical AIMD and path-integral AIMD
trajectories. The DFT engine for these reference data is not stated
in the paper (it inherits the published Cheng et al.\ dataset).
\lstinputlisting[language=turtle]{artifacts/kg/listings/batzner2022/q04-reference-settings.ttl}

\paragraph{Q5 -- Training dataset.}
The water + ices joint training set used in the data-efficiency
experiment consists of 133 structures sampled uniformly from a
140{,}000-frame full set (100{,}000 liquid water; 20{,}000 ice
Ih~b; 10{,}000 ice Ih~c; 10{,}000 ice Ih~d). The training labels
are energies and forces.
\lstinputlisting[language=turtle]{artifacts/kg/listings/batzner2022/q05-dataset.ttl}

\paragraph{Q6 -- Sampling strategies.}
A single sampling strategy is in play: vibrational sampling via
classical and path-integral AIMD trajectories at multiple
temperatures and pressures.
\lstinputlisting[language=turtle]{artifacts/kg/listings/batzner2022/q06-sampling.ttl}

\paragraph{Q7 -- MLIP method, components, implementation.}
NequIP is an $E(3)$-equivariant message-passing GNN that uses
tensor-product convolutions over geometric tensors of irreducible
$O(3)$ representations of order $\ell=0\dots l_{\max}$, parities
$p\in\{-1,1\}$, with SiLU/tanh gate non-linearities, a ResNet-style
self-interaction--convolution--concatenation update, and a final
atom-wise output block. The training algorithm is the AMSGrad
variant of Adam ($\beta_1=0.9$, $\beta_2=0.999$, $\epsilon=10^{-8}$,
no weight decay) with an on-plateau learning-rate scheduler
(patience 50, decay factor 0.8). The loss is a weighted MSE on
energies and force components ($\lambda_F=100\,000$, $\lambda_E=1$
for the model~c run reported below). Implementation is the
\texttt{nequip} PyTorch code, version 0.3.3.
\lstinputlisting[language=turtle]{artifacts/kg/listings/batzner2022/q07-method.ttl}

\paragraph{Q8 -- Hyperparameter settings.}
The water + ices model uses 6 interaction blocks, $l_{\max}=2$, 32
features (with both even and odd parity), and a 6~\AA{} radial
cutoff. Radial features are 8 trainable Bessel basis functions
processed by a 3-hidden-layer MLP of size 64 with SiLU
non-linearities. Learning rate 0.005, batch size 1 (with weight
0.99 EMA on training weights for evaluation).
\lstinputlisting[language=turtle]{artifacts/kg/listings/batzner2022/q08-hyperparameter-settings.ttl}

\paragraph{Q9 -- MLIP run and trained model.}
A single training run on 133 structures produces one trained
NequIP water+ices model. Three loss-weighting variants are reported
(model~a: $\lambda_F=1,\lambda_E=0$; model~b: $\lambda_F=100$;
model~c: $\lambda_F=100\,000$); we encode model~c as the headline.
\lstinputlisting[language=turtle]{artifacts/kg/listings/batzner2022/q09-run-and-model.ttl}

\paragraph{Q10 -- Benchmark results.}
Headline figures for the water + ices test set (model~c,
$\lambda_F=100\,000$, $\lambda_E=1$, 133 train) are RMSEs of
1.7~meV per H$_2$O molecule on energies and 12.2~meV/\AA{} on
liquid-water force components. NequIP is competitive with DeepMD
trained on $\sim$1000$\times$ more data and significantly
outperforms it on three of four phases. Per-phase results for the
three ice configurations are an additional Q12 gap.
\lstinputlisting[language=turtle]{artifacts/kg/listings/batzner2022/q10-benchmark-results.ttl}

\paragraph{Q11 -- Computational resources.}
All NequIP models were trained on an NVIDIA Tesla V100 GPU in
single-GPU training. The paper does not report a specific training
duration or GPU-hour count for the water+ices run (only that
``competitive results can typically be obtained within a matter of
hours or often even minutes''), nor peak memory or inference time
per atom for the headline pass.
\lstinputlisting[language=turtle]{artifacts/kg/listings/batzner2022/q11-resources.ttl}

\paragraph{Q12 -- Gaps.}
Beyond Q11, the following ontology-expressible items are not
represented in this encoding pass: (i) the six additional
benchmarks reported in the same paper (MD17/rMD17 small molecules,
Molecules@CCSD/CCSD(T), formate decomposition on Cu$\langle 110
\rangle$, Li$_4$P$_2$O$_7$ amorphous glass, LiPS superionic
conductor); (ii) per-phase benchmark results for the three ice Ih
configurations and the data-efficiency learning-curve sweep across
$\{10, 100, 1000, 2500\}$ training-set sizes for LiPS; (iii) the
DFT engine for the water reference data ($\dlRole{usedDFTCode}$);
(iv) detailed PBE0-TS reference settings (k-point mesh, energy
cutoff, pseudopotential type); (v) the PyTorch / PyTorch
Geometric / e3nn / Python versions used for training (the paper
does report all of these but our protocol does not capture
secondary toolchain libraries beyond the MLIP \texttt{Library}
hook); (vi) the AMSGrad-Adam learning-rate schedule on a per-system
basis. The paper also reports a ``temperature transferability''
benchmark on 3BPA but only as cited in
the Allegro paper (Musaelian et al.~2023); this is not encoded here.
